% Dataset section. Paste into your own thesis, or \input it.

\section{Dominance and anatomical variants}
\label{sec:dataset-dominance}

The ImageCAS-X cohort is annotated as 729 right-dominant, 41 left-dominant and 30 codominant
cases, or 91.1\%, 5.1\% and 3.8\% of the 800 usable cases \cite{bransby2026imagecasx}.
The underlying scans were acquired at a single hospital in Guangdong, China
\cite{zeng2023imagecas}, and the observed proportion of right dominance falls between the
European and East Asian figures reported by Wu et al.\ (Section~\ref{sec:dominance}).

Dominance is recorded twice in ImageCAS-X, independently, and the two records agree
\cite{bransby2026imagecasx}: it is an explicit annotation, and it is implied by which arteries are
labeled, since the posterior vessels arise from whichever system is dominant.
In Table~\ref{tab:posterior} the right posterior vessels appear in essentially every
right-dominant case and in no left-dominant case, the left posterior vessels invert the pattern
exactly, and codominant cases draw on both.

\begin{table}[h]
  \centering
  \caption{The artery labels and the dominance annotation are two independent records of the same
    anatomical fact. Prevalence of the posterior vessels by annotated dominance over the 800
    usable cases, from \texttt{scripts/cohort\_numbers.py}; the right and left posterior vessels
    are mutually exclusive in the single-dominance groups.}
  \label{tab:posterior}
  \begin{tabular}{lrrrr}
    \hline
    Vessel & All (\%) & Right (\%) & Left (\%) & Co (\%) \\
    \hline
    R-PDA & 93.0 & 99.7 &  0.0 & 56.7 \\
    R-PLA & 92.0 & 99.2 &  0.0 & 43.3 \\
    L-PDA &  5.8 &  0.0 & 90.2 & 30.0 \\
    L-PLA &  8.8 &  1.0 & 90.2 & 86.7 \\
    \hline
  \end{tabular}
\end{table}

\begin{figure}[h]
  \centering
  \includegraphics[width=\textwidth]{content/background/figures/05_dominance_comparison.png}
  \caption{The posterior vessels follow the dominant system, so dominance changes which arteries
    exist rather than how large they are. Right-, left- and codominant coronary trees from the
    cohort; the right posterior descending artery is absent from the left-dominant tree entirely.}
  \label{fig:dominance}
\end{figure}

\subsection{Absent left main}

The absent left main variant (Section~\ref{sec:variants}) is present in 11 of the 800 cases, or
1.4\%, and it is identifiable from two independent signals that agree exactly: those cases carry
no LM label, and their left centerlines resolve into two components each with its own ostium
(\texttt{scripts/survey\_topology.py}).

\subsection{Ramus intermedius}

An IM label (Section~\ref{sec:variants}) is present in 196 of the 800 ImageCAS-X cases, or
24.5\%, consistent with the pooled CCTA prevalence of 26.5\% reported by Dai et al.\
\cite{dai2025ramus}.
The cohort figure is also stable across dominance groups, at 24.7\%, 22.0\% and 23.3\% for
right-, left- and codominant cases respectively.
